With the three structural rules (at least two qubits, at least one two-qubit gate, no consecutive same-axis rotations on a qubit) the model scores $0.939 \pm 0.015$ on 3 vs 8 and $0.650 \pm 0.067$ on ten classes, indistinguishable from the unconstrained $0.924$ and $0.653$, and all ten circuits are valid on the first attempt. What it built is the more interesting result. In every binary trial the model kept the linear accumulator and found a circuit identity that satisfies the rules: it interleaves $z$ gates between the $ry$ rotations (since $z\,ry(b)\,z = ry(-b)$, the chain still sums, with alternating signs absorbed into the weights; ``alternating Z gates and signed RY angles accumulate this evidence into one decision angle'', it explains), or accumulates the same weighted sum as a phase with $p$ or $rz$ gates between CNOTs and converts it to a probability with a final Hadamard, and then moves the result to the readout qubit with a single CNOT. For ten classes two of the five circuits are of the same kind (``the $\pi$-angle ZZ gates provide sign refocusing between successive encoding rotations''), and three use genuinely multi-qubit designs: 117 $u3$ rotations with 12 CNOTs on five qubits ($0.706$), and two eight-qubit circuits in which ``four address qubits select a class-specific affine phase; interference converts that phase into a score monotonic in a sparse nearest-prototype discriminant'', built from 39 Toffoli and 146 controlled-phase gates ($0.713$, the best ten-class circuit of the study). The model, in other words, knows exactly which weight space it wants to program---a signed sum of pixels---and, forbidden the obvious circuit for it, writes a different circuit that computes the same thing, or, when it must, a nearest-prototype classifier whose prototypes it writes by hand.
